\itemhead{习题1.2 · 第7题（教材p.24）}{XJ-B1-C01-HW-P024-Y0897}
\label{XJ-B1-C01-HW-P024-Y0897}
{\footnotesize 湘教版教材p.24 / PDF第31页 / 习题1.2 / 第7题}
\field{原题}{下列各题中，$p$ 是 $q$ 的什么条件（在“充分而不必要条件”“必要而不充分条件”“充要条件”“既不充分又不必要条件”中选出一种）？为什么？
\subq{1}已知集合 $A=\{1,a\}$，$B=\{1,2,3\}$，$p$：$a=3$，$q$：$A\subseteq B$；
\subq{2}$p$：$a\in\mathbb N_+$，$q$：$a\in\mathbb Z_+$；
\subq{3}设 $x,y$ 是实数，$p$：$x>y$，$q$：$|x|>|y|$；
\subq{4}$p$：$D$ 在 $\triangle ABC$ 的边 $BC$ 的中线上，$q$：$\triangle ABD$ 的面积 $=\triangle ACD$ 的面积。}
\field{参考答案}{依次为充分而不必要、充要、既不充分又不必要、充分而不必要条件。（4）中线上的点必使两面积相等，但中线所在直线上三角形外的点也可使面积相等。}
{\small \field{关联知识}{L02 充分、必要、充要条件（教材帮p.30、31 / PDF39、40）}
\field{方法与易错点}{分别检查两个方向；用解集包含或具体反例判定，不凭句子位置猜方向。 易错：把“必要”箭头反写；只有一个方向证据便断言“不必要”；忽略给定论域。}
\field{答案依据}{答案册相关题组：PDF第7页（册内页6）；按当前题库原题逐项重算，不按旧版题号机械对套。}}
